% !TEX root = ../main.tex
\chapter{Conclusions and Future Work}
\label{sec:conclusions}

\section{Conclusions}
\label{sec:conclusions-conclusions}
This thesis asked how training an MoE language model with Muon instead of AdamW changes the work of the bias that keeps its experts equally busy. We measured this work by the span, the difference between the largest and the smallest bias in a layer. We trained 145 runs of models with 55M and 108M active parameters on educational web text. In each run, we assigned an optimizer to each part of the model, and we call this assignment a partition. We compared the runs seed by seed.

RQ1 compared Muon with AdamW\@. Under Muon, unlike under AdamW, the update of an expert did not grow with its load. We expected this to weaken the tendency of busy experts to become busier, so that the bias would work less. It worked harder instead. The span was higher at every seed: by about a third averaged over the layers, and more than twice as large in the first MoE layer. We call this higher span of the first MoE layer the rise. Muon also reached a lower loss and made the experts more specialized. Most of the rise came with the learning rate of Muon, which our tuning had set to twice that of AdamW\@. At half this rate, which trained slightly better, about three quarters of the rise disappeared.

RQ2 trained AdamW, Muon, and Muon with the router on AdamW, the three partitions of current models, at up to 16 times their tuned learning rates. No run diverged. The two Muon partitions reached the same loss at every multiple, and at high multiples their loss exceeded that of AdamW\@. The Qwen3.8-Flash-Next report moved the router to AdamW, saying that Muon made it unstable early in training. In our runs, however, the loads shifted more in the warmup with the router on AdamW than under Muon.

RQ3 moved each part of the model to Muon in turn. The optimizer of the router changed the span of the first MoE layer little. What mattered was the optimizer of the backbone, that is, of every weight matrix except the routers, the embedding table, and the output head. With the whole backbone on Muon, the inputs of the first router shared a direction. Through this direction, the router gave the same experts a higher score for every token. The bias needed hundreds of steps to cancel this preference. Meanwhile, the busiest expert of this layer received four to seven times as many tokens as it would if the loads were even. In an exploratory test of this explanation, we subtracted the mean input of the batch from the input of the first router, which removes the shared direction. Over the first 3,000 steps, this centering eliminated the overload and most of the rise, without changing the loss. Muon on the routed experts alone raised the span even more, without a strong shared direction, and we do not know how.

The rise persisted with a bias four times faster and with the balance loss of DeepSeek-V3, although both reduced it on average. On the large model, the rise was larger in absolute terms. Exploratory analyses also found the first MoE layer still overloaded at the end of training.

With the backbone on Muon, the update of a router on AdamW acted more along the shared direction than that of a router on Muon, on both model sizes. This fits the larger shifts of the loads during the warmup and, on the small model, the higher overload of the first MoE layer. The runs with the router on AdamW also reached a slightly higher loss at every seed, for a reason that our results do not identify. For a recipe that trains the backbone with Muon and balances the loads with the bias, our results therefore favor keeping the router on Muon. They also point to the first MoE layer as the one to watch most closely. These conclusions rest on models far smaller than current ones (\cref{sec:discussion-limitations}).

\section{Future Work}
\label{sec:conclusions-future}
Our results leave several questions open, and each suggests further experiments.
\begin{itemize}
  \item \textbf{Centering over whole runs.} We tested the centering only over the first 3,000 steps and in one layer. The next test is to center the input of every router over whole runs, at both sizes. The mean over the batch, however, makes the routing of a token depend on its batch. A practical version would subtract a running mean instead.
  \item \textbf{Other balancers.} The bias four times faster lowered the overload of the first MoE layer, but on the large model the loads kept shifting (\cref{sec:discussion-overload}). A step of the bias that is large early in training and small later may lower the overload without this shifting. Such a step should be compared with the rules of Kimi K3 and Trinity Large (\cref{sec:related-balancer}), and with a larger weight of the balance loss. Such a weight may suppress the shared direction.
  \item \textbf{Larger models and longer runs.} The shared direction and the overload of the first MoE layer grew from the small model to the large one. Larger and longer runs of the three partitions would show whether this overload keeps growing or settles. These runs should tune their own learning rates, since a lower rate of Muon removed most of the rise on the small model. They may also have enough loss spikes to relate them to the routing.
  \item \textbf{Other routers, architectures, and data.} A router with unbounded scores, as in DeepSeek-V4, would show whether the span grows beyond the ceiling of our sigmoid router (\cref{sec:discussion-limitations}). In a softmax router, large logits cause rounding errors (\cref{sec:discussion-router}). Such a router would show whether Muon also gives it large logits at high learning rates. Models with no dense block, or with several, before the first MoE layer would tell us how this choice shapes the shared direction. Other data, such as code, may also change the shared direction.
  \item \textbf{The open parts of the mechanism.} We do not know why only the whole backbone on Muon gives the inputs of the first router a shared direction. We also do not know how Muon on the routed experts alone raises the span even more. The control is an optimizer that gives the update of every expert the same size, as Muon does, but without the orthogonalization. Its tuning runs suggest that this equal size plays a role (\cref{sec:discussion-mechanism}), and training the routed experts alone with the control would test it.
  \item \textbf{The cost of the overload.} In our models, which drop no tokens, removing the overload did not change the loss. Two other costs remain to be measured (\cref{sec:discussion-overload}). When the experts sit on different devices, the busiest ones slow down training. When each expert has a capacity, a limit on its tokens, the extra tokens are dropped.
\end{itemize}
